\documentclass[10pt,a4paper]{amsart}
\usepackage[margin=19mm]{geometry}
\usepackage{amsmath,amssymb,mathtools}
\usepackage[scr=rsfs,cal=euler]{mathalfa}
\usepackage{xfrac}
\usepackage[unicode,hidelinks,bookmarks=false]{hyperref}
\newcommand{\ie}{i.e.\ }
\newcommand{\cf}{cf.\ }
\newcommand{\resp}{respectively\ }
\newcommand{\Subsection}{\subsection}
\providecommand{\nobreakdash}{\nobreak\hbox{-}}
\setlength{\emergencystretch}{2em}
\multlinegap0pt
\usepackage{bm}
\usepackage{bbm}
\usepackage{dsfont}
\usepackage{xspace}
\usepackage{cancel}
\usepackage{mathtools}
\usepackage{cleveref}
\usepackage[shortlabels]{enumitem}
\usepackage{amsthm}
\makeatletter
\@ifundefined{theorem}{\newtheorem{theorem}{Theorem}}{}
\@ifundefined{definition}{\newtheorem{definition}[theorem]{Definition}}{}
\@ifundefined{lemma}{\newtheorem{lemma}[theorem]{Lemma}}{}
\makeatother
\DeclareMathOperator*{\E}{\mathbb{E}}
\newcommand{\C}{\mathbb C}
\newcommand{\R}{{\mathbb R}}
\newcommand{\Tr}{\operatorname{Tr}}
\newcommand{\cH}{\mathcal H}
\newcommand{\cR}{\mathcal R}
\newcommand{\sfx}{\mathsf x}
\newcommand{\spn}{\operatorname{span}}
\newcommand{\tr}{\mathrm{tr}}
\renewcommand{\geq}{\geqslant}
\renewcommand{\leq}{\leqslant}
\title[Sharp Bounds on Ground State Energy of the SYK Model]{Sharp Bounds on Ground State Energy of the SYK Model}
\author{Arpon Basu, Pravesh K. Kothari, Siddhant Midha}
\date{Source: arXiv:2607.27185v1 (CC BY 4.0)}
\begin{document}
\maketitle

\noindent\emph{Source and licence.} Sharp Bounds on Ground State Energy of the SYK Model. Version arXiv:2607.27185v1 (CC BY 4.0), distributed under the Creative Commons Attribution 4.0 International licence, as recorded in the arXiv licence metadata for this exact version and shown on the abstract page https://arxiv.org/abs/2607.27185v1. Full rights notice: licenses/arxiv-2607.27185-CC-BY-4.0.txt.\par\smallskip

\noindent\textit{Editorial scope.} Counted unit: the Theorem whose source label is \texttt{thm:syk-spectral-norm-lower-random} in the article (source0.tex lines 1964--1975, source label \texttt{thm:syk-spectral-norm-lower-random}), delivered together with the article's own complete original proof of it (source0.tex lines 1976--2064, one \texttt{proof} environment of 89 lines). No mathematics is written, repaired or re-derived: statement and proof are the article's own text line for line. Required context retained uncounted: eq:syk-def (lines 830--832); lem:random-proj-compression (lines 867--870); defn:relevant-subspace (lines 1636--1642); lem:x-large-eig (lines 1644--1647); lem:lower-bound-syk-general (lines 1754--1757); eq:complete-sign-eig (lines 1793--1795); lem:complete-matrix-values (lines 1801--1814); thm:syk-spectral-norm-random (lines 1906--1912). Every definition, display and label of those lines is retained unchanged. Omitted from this package and not redistributed: 1--829, 833--866, 871--1635, 1643--1643, 1648--1753, 1758--1792, 1796--1800, 1815--1905, 1913--1963, 2065--2087. Nothing inside a retained excerpt is omitted or altered: every retained body is delivered exactly as the article's lines are written, so no omission receipt of any kind accompanies this package. Citations: the retained excerpts carry no citation command at all, so no bibliography entry of the article is transcribed and no reference is redistributed. Reference adaptation: none was needed and none was made; every reference inside the delivered unit resolves inside it, so no unresolved reference or citation is produced. Printed slips kept literal: no printed word, symbol, hypothesis or number of the counted statement or of its proof was altered. Layout and class: the article's own class is \texttt{article} with packages including nag, dtrt, inputenc, stmaryrd, xspace, enumerate, fontenc, textcomp, babel, mathtools, amsthm, thmtools, tikz, pdflscape; the wrapper is the lane's pinned \texttt{amsart} editorial wrapper at 10pt on A4, which restates the theorem-like environments the retained text needs and adds the article's own macro definitions that the retained text needs (\texttt{\textbackslash C}, \texttt{\textbackslash E}, \texttt{\textbackslash R}, \texttt{\textbackslash Tr}, \texttt{\textbackslash cH}, \texttt{\textbackslash cR}, \texttt{\textbackslash geq}, \texttt{\textbackslash leq}, \texttt{\textbackslash sfx}, \texttt{\textbackslash spn}, \texttt{\textbackslash tr}), with extra wrapper packages (bm, bbm, dsfont, xspace, cancel, mathtools, cleveref, enumitem). The delivered numbering is the unit's own. Assets: no retained span contains a figure environment, a graphics-inclusion command, a PDF-page inclusion, a table or a TikZ picture; no image file is copied into this package, the package declares no asset dependency and no image file is redistributed with it. Licence: version arXiv:2607.27185v1 of this article is distributed under the Creative Commons Attribution 4.0 International licence, as recorded in the arXiv licence metadata for this exact version and shown on the abstract page https://arxiv.org/abs/2607.27185v1. A full rights notice is delivered at licenses/arxiv-2607.27185-CC-BY-4.0.txt. Characters: every retained excerpt is pure ASCII, so the delivered engine, XeLaTeX with its default Latin Modern text font, reports no missing character.
\begin{equation}\label{eq:syk-def}
    H^\cH_{\operatorname{SYK}}:= \frac{1}{\sqrt{|\cH|}}\sum_{S\in\cH}g_S\Gamma_S.
\end{equation}

\begin{lemma}\label{lem:random-proj-compression}
    Let $G\in\C^{\Omega\times\Omega}$ be a non-zero Hermitian PSD matrix with $\|G\|_{\operatorname{op}} = \mu$. Further suppose that $G$ has a constant diagonal, i.e. $\langle e_\omega, Ge_\omega\rangle = \tr(G)$ for every $\omega\in\Omega$. Then for $A_G:= (D/m)P^*GP\in\C^{\Omega'\times\Omega'}$ we have with probability $\geq 1 - O(\exp(-100\log(2\Tr(G)/\mu)))$ that
    \[\|A_G\|_{\operatorname{op}}\leq O(1)\cdot\max\left\{\mu, \frac{\Tr(G)}{m}\cdot\log\left(\frac{2\Tr(G)}{\mu}\right)\right\}\]
\end{lemma}

 \begin{definition}[Krylov Space]\label{defn:relevant-subspace}
 For parameters $t_0, M$ define the subspace (on the twisted Fock space)
\begin{equation}
    \cR_{t_0, M}:= \spn\{v_t: t_0\leq t < t_0 + M\}
\end{equation}
and $\Pi_{\cR_{t_0, M}}$ to be the orthogonal projector onto $\cR_{t_0, M}$. If $t_0, M$ are understood from context we will not mention it explicitly.
 \end{definition}

\begin{lemma}[Spectral edge of hopping model]\label{lem:x-large-eig}
     Suppose $B\geq\kappa\geq\Omega_n(1)$ for some parameter $\kappa$. Write $t_0:= \lceil \kappa/\nu\rceil, L:= t_0 + M$. Then 
    \[\lambda_{\max}\left(\Pi_\cR \sfx \Pi_\cR\right)\geq\frac{2}{\sqrt{\nu}}\cdot\left(1 - O(1)\cdot\max\left\{\kappa^{-1}, \nu^2\right\}\right).\]
\end{lemma}

\begin{lemma}[Lower Bounds on Operator Norm of the SYK Hamiltonian]\label{lem:lower-bound-syk-general}
    Let $n\geq k\geq 2$ be even integers, and let $\alpha:= k^2/n < 1/16$. Consider the SYK Hamiltonian $H^\cH_{\operatorname{SYK}}$ (from \cref{eq:syk-def}) for the hypergraph $\cH$. Let $\cR$ be any space of the kind defined in \cref{defn:relevant-subspace}. Then
    \[\E\|H^\cH_{\operatorname{SYK}}\|_{\operatorname{op}}\geq \lambda_{\max}(\Pi_\cR\sfx\Pi_\cR) - O(1)\cdot \sqrt{\frac{L(\rho^+ + \delta)}{1 - q}}.\]
\end{lemma}

\begin{equation}\label{eq:complete-sign-eig}
    \lambda_j:= \binom{n}{k}^{-1}\cdot\sum_{t = j}^k(-2)^t\binom{k - j}{t - j}\cdot\binom{n - t - j}{k - t}
\end{equation}

\begin{lemma}[Sign Matrix Eigenvalues]\label{lem:complete-matrix-values}
    Suppose $n\geq k\geq 2$. Write $\alpha:= k^2/n$, and assume $\alpha < 1/16$. For $\lambda_j$ as in \cref{eq:complete-sign-eig}, we have
    \begin{equation}\label{eq:lambdaj-estimate}
        \left|\frac{\lambda_j}{(-2)^j\binom{n - 2j}{k - j}\binom{n}{k}^{-1}} - 1\right|\leq 8\alpha.
    \end{equation}
    In particular, $\lambda_j > 0$ if and only if $j$ is even. Moreover, 
    \begin{enumerate}[(1)]
        \item\label{item:lambda1-bound} $(1 - k/n)(1 - 8\alpha)\leq |\lambda_1|/(2k/n)\leq 1 + 8\alpha$.
        \item\label{item:lambdaj-bound} $|\lambda_j|\leq 2(8k/n)^j$ for all $j\geq 0$. Consequently, $\sum_{j = 0}^{k}d_j|\lambda_j|\leq e^{9k}$, where $d_j:= \binom{n}{j} - \binom{n}{j - 1}$.
        \item\label{item:even-lambdaj-upperbound} $\max_{\substack{2\leq j\leq k\\j\text{ even}}}\lambda_j = \Theta(\alpha/n)$.
        \item\label{item:odd-lambdaj-upperbound} $\max_{\substack{1\leq j\leq k\\j\text{ odd}}}|\lambda_j| = |\lambda_1| \leq (2 + O(\alpha))\cdot\sqrt{\alpha/n}$.
        \item $2\alpha - 2\alpha^2\leq 1 - \lambda_0\leq 2\alpha$. 
    \end{enumerate}
\end{lemma}

\begin{theorem}[Operator Norm of the Sparse SYK: Upper Bound]\label{thm:syk-spectral-norm-random}
    Let $n\geq k\geq 2$ be even integers, and suppose $\alpha:= k^2/n < 1/16$.  Let $\cH$ be a random $k$-uniform hypergraph with $m\geq 2^{Ck}n\log n$ edges, where $C > 0$ is some sufficiently large absolute constant. Let 
    \[H^\cH_{\operatorname{SYK}}:= \frac{1}{\sqrt{|\cH|}}\sum_{S\in\cH}g_S\Gamma_S\]
    be the SYK Hamiltonian corresponding to $\cH$ (as in \cref{eq:syk-def}). Then with probability $\geq 1 - O(1/\log(n))$ over the randomness of $\cH$ we have 
    \[\E_{\bm g}\|H^\cH_{\operatorname{SYK}}\|_{\operatorname{op}}\leq \sqrt{\frac{2}{\alpha}}\cdot\left(1 + O\left(\sqrt{\alpha + e^{-\Omega(k)}}\right)\right) = \frac{\sqrt{2n}}{k}\cdot\left(1 + O\left(\sqrt{\frac{k^2}{n} + e^{-\Omega(k)}}\right)\right).\]
    In particular if $k\geq\omega(\log(n))$ then $\E\|H^\cH_{\operatorname{SYK}}\|_{\operatorname{op}}\leq \sqrt{2n}/k + O(1)$.
\end{theorem}

\begin{theorem}[Operator Norm of the Sparse SYK: Lower Bound]\label{thm:syk-spectral-norm-lower-random}
    Let $n\geq k\geq 2$ be even integers, and let $\alpha:= k^2/n < 1/16$. Consider the SYK Hamiltonian (from \cref{eq:syk-def}) for a random hypergraph $\cH$ with $m\geq 2^{Ck}n\log(n)$ hyperedges for some large enough absolute constant $C > 0$. Then 
    \[\E\|H\|_{\operatorname{op}}\geq
\sqrt{\frac2\alpha}
\left(
1-O(1)\max\left\{
k^{-1/4},
\alpha^2
\right\}
\right).\]
with probability $\geq 1-O(1/\log n)$ over the randomness of $\cH$.
\end{theorem}

\begin{proof}
Throughout the proof we'll borrow notation from \cref{lem:x-large-eig}. By \cref{lem:lower-bound-syk-general} we know that if $B\geq\Omega(1)$, then
\begin{equation}
\label{eq:random-lower-transfer}
\E_{\bm g}\left\|H^\cH_{\operatorname{SYK}}\right\|_{\operatorname{op}}\geq\lambda_{\max}\!\left(\Pi_\cR\sfx\Pi_\cR\right)-O(1)\sqrt{\frac{L(\rho^++\delta)}{1-q}}.
\end{equation}

Note that for any unit vector $v\in\R^\cH$ such that $\langle v, u\rangle = 0$, we have
\begin{equation*}
-\langle v,\mathcal E^\cH v\rangle = \langle v,A^-v\rangle-\langle v,A^+v\rangle \leq \langle v,A^-v\rangle \leq \|A^-\|_{\operatorname{op}}\implies\rho^-\leq \|A^-\|_{\operatorname{op}}.
\end{equation*}

We now apply \cref{lem:random-proj-compression} to $G^-$. By
\cref{lem:complete-matrix-values},
\[\|G^-\|_{\operatorname{op}} = \max_{\substack{j\geq1\\\lambda_j<0}}|\lambda_j| = |\lambda_1| \leq O\left(\sqrt{\frac{\alpha}{n}}\right) = O\left(\frac{k}{n}\right),\;\Tr(G^-) = \sum_{\substack{j\geq1\\\lambda_j<0}}|\lambda_j|\dim(V_j)\leq e^{9k}.\]
Moreover, by the $\mathfrak{S}_n$-invariance of the Johnson eigenspaces,
the diagonal entries of $G^-$ are all equal. Hence
\cref{lem:random-proj-compression} gives, with probability at least
$1-O(n^{-100})$,
\[\|A^-\|_{\operatorname{op}}\leq O(1)\max\left\{\frac{k}{n}, \frac{e^{9k}}{m}\log\left(e^{O(k)}n^2\right)\right\}\leq O(1)\max\left\{\frac{k}{n}, \frac{e^{O(k)}(k+\log n)}{m}\right\}\leq O\left(\frac{k}{n}\right).\]
where the last inequality follows from
$m\geq2^{Ck}n\log n$ after choosing $C$ sufficiently large. Consequently we have
\[\rho^-\leq O\left(\frac{k}{n}\right) = O\left(\sqrt{\frac{\alpha}{n}}\right).\]

As proved in \cref{thm:syk-spectral-norm-random}, with probability at
least $1-O(1/\log n)$,
\[\delta\leq O\left(\sqrt{\frac{\alpha\log n}{m}}\right)\leq O\left(\frac{e^{-\Omega(k)}}{n}\right),\qquad \rho^++\delta\leq O(1)\cdot\frac{\alpha+e^{-\Omega(k)}}{n}\leq O\left(\frac{k}{n}\right),\]
\begin{align*}
(2 + o_n(1))\alpha\geq 2\alpha + O\!\left(
\frac{e^{-\Omega(k)}}{\sqrt{n^3\log n}}
\right)\geq\nu = 1-q
&\geq
2(\alpha-\alpha^2)
-
O\!\left(
\frac{e^{-\Omega(k)}}{\sqrt{n^3\log n}}
\right)
\geq \Omega(\alpha).
\end{align*}
Consequently we have $\nu = \Theta(\alpha)$. Thus
\[
B
:=
\sqrt{\frac{\nu}{\rho^- + \delta}}
 \geq
\Omega\!\left(
\sqrt{\frac{\alpha}{k/n}}
\right)
\geq
\Omega(\sqrt{k})
\geq\Omega(1).
\]
Hence the hypothesis needed by \cref{lem:x-large-eig} is satisfied. Now choose $\kappa = \Theta(\sqrt{k})\leq B$, and note that this choice yields 
\[L\leq O\left(\frac{\kappa}{\nu}\right)\leq O\!\left(\frac{n}{k^{3/2}}\right).\]

Thus
\begin{align*}
\frac{L(\rho^++\delta)}{1-q}
&\leq
O\!\left(\frac{n}{k^{3/2}}\right)
\frac{(\alpha+e^{-\Omega(k)})/n}{\alpha}
 =
O\!\left(
\frac{\alpha+e^{-\Omega(k)}}{\alpha k^{3/2}}
\right)
\leq
O\!\left(\frac1{\alpha\sqrt{k}}\right),
\end{align*}
where in the last inequality we used
$\alpha+e^{-\Omega(k)}\leq O(k)$ for $k\geq2$. Therefore
\begin{equation}
\label{eq:random-transfer-error-root}
\sqrt{\frac{L(\rho^++\delta)}{1-q}}\leq O\!\left(\frac1{\sqrt{\alpha}\,k^{1/4}}\right) = \sqrt{\frac2\alpha}\,O\!\left(k^{-1/4}\right).
\end{equation}

Finally \cref{lem:x-large-eig} applied with the preceding bounds on $\nu$, $\rho^-+\delta$, $\kappa$, and $L$, gives
\begin{equation}
\label{eq:random-x-large-eigenvalue}
\lambda_{\max}\!\left(\Pi_\cR\sfx\Pi_\cR\right)\geq\sqrt{\frac2\alpha}\left(1-O(1)\max\left\{k^{-1/2},\alpha^2\right\}\right).
\end{equation}
Combining
\eqref{eq:random-lower-transfer},
\eqref{eq:random-transfer-error-root}, and
\eqref{eq:random-x-large-eigenvalue}, and absorbing the transfer error into
the existing $O(k^{-1/4})$ relative error, yields
\[\E_{\bm g}\left\|H^\cH_{\operatorname{SYK}}\right\|_{\operatorname{op}}\geq\sqrt{\frac2\alpha}\left(1-O(1)\max\left\{k^{-1/4},\alpha^2\right\}\right).\]
All the required inequalities hold simultaneously with probability
at least $1-O(1/\log n)$, and thus we are done.
\end{proof}

\end{document}
